\documentclass[10pt,a4paper]{amsart}
\usepackage[margin=19mm]{geometry}
\usepackage{amsmath,amssymb,mathtools}
\usepackage[scr=rsfs,cal=euler]{mathalfa}
\usepackage{xfrac}
\usepackage[unicode,hidelinks,bookmarks=false]{hyperref}
\newcommand{\ie}{i.e.\ }
\newcommand{\cf}{cf.\ }
\newcommand{\resp}{respectively\ }
\newcommand{\Subsection}{\subsection}
\providecommand{\nobreakdash}{\nobreak\hbox{-}}
\setlength{\emergencystretch}{2em}
\multlinegap0pt
\usepackage{bm}
\usepackage{bbm}
\usepackage{dsfont}
\usepackage{xspace}
\usepackage{cancel}
\usepackage{mathtools}
\usepackage{amsthm}
\makeatletter
\@ifundefined{theorem}{\newtheorem{theorem}{Theorem}}{}
\@ifundefined{lemma}{\newtheorem{lemma}[theorem]{Lemma}}{}
\@ifundefined{proposition}{\newtheorem{proposition}[theorem]{Proposition}}{}
\makeatother
\title[Combinatorics of descent algebras and graph coverings]{Combinatorics of descent algebras and graph coverings}
\author{Philippe Biane}
\date{Source: arXiv:2506.05528v1 (CC BY 4.0)}
\begin{document}
\maketitle

\noindent\emph{Source and licence.} Combinatorics of descent algebras and graph coverings. Version arXiv:2506.05528v1 (CC BY 4.0), distributed under the Creative Commons Attribution 4.0 International licence, as recorded in the arXiv licence metadata for this exact version and shown on the abstract page https://arxiv.org/abs/2506.05528v1. Full rights notice: licenses/arxiv-2506.05528-CC-BY-4.0.txt.\par\smallskip

\noindent\textit{Editorial scope.} Counted unit: the Proposition whose source label is \texttt{recoil} in the article (source0.tex lines 116--126, source label \texttt{recoil}), delivered together with the article's own complete original proof of it (source0.tex lines 127--147, one \texttt{proof} environment of 21 lines). No mathematics is written, repaired or re-derived: statement and proof are the article's own text line for line. Required context retained uncounted: >2 (lines 101--106). Every definition, display and label of those lines is retained unchanged. Omitted from this package and not redistributed: 1--100, 107--115, 148--877. Nothing inside a retained excerpt is omitted or altered: every retained body is delivered exactly as the article's lines are written, so no omission receipt of any kind accompanies this package. Citations: the retained excerpts carry no citation command at all, so no bibliography entry of the article is transcribed and no reference is redistributed. Reference adaptation: none was needed and none was made; every reference inside the delivered unit resolves inside it, so no unresolved reference or citation is produced. Printed slips kept literal: no printed word, symbol, hypothesis or number of the counted statement or of its proof was altered. Layout and class: the article's own class is \texttt{amsart} with packages including inputenc, fontenc, amssymb, amsthm, amsmath, tikz, multicol, wrapfig, url, hyperref, geometry, pdfsync; the wrapper is the lane's pinned \texttt{amsart} editorial wrapper at 10pt on A4, which restates the theorem-like environments the retained text needs and adds the article's own macro definitions that the retained text needs (none), with extra wrapper packages (bm, bbm, dsfont, xspace, cancel, mathtools). The delivered numbering is the unit's own. Assets: no retained span contains a figure environment, a graphics-inclusion command, a PDF-page inclusion, a table or a TikZ picture; no image file is copied into this package, the package declares no asset dependency and no image file is redistributed with it. Licence: version arXiv:2506.05528v1 of this article is distributed under the Creative Commons Attribution 4.0 International licence, as recorded in the arXiv licence metadata for this exact version and shown on the abstract page https://arxiv.org/abs/2506.05528v1. A full rights notice is delivered at licenses/arxiv-2506.05528-CC-BY-4.0.txt. Characters: every retained excerpt is pure ASCII, so the delivered engine, XeLaTeX with its default Latin Modern text font, reports no missing character.
\begin{lemma}\label{>2}
Let $\sigma\in S_n$ and  $i\in\{1,\ldots,n-1\}$ then 
$R(\sigma)=R(\sigma s_i)$
 if and only if
$|\sigma_i-\sigma_{i+1}|\geq 2$.
\end{lemma}

\begin{proposition}\label{recoil}
Let $\sigma,\pi,\rho\in S_n$ and $i\leq n-1$ be such that $$\sigma=\pi\rho \quad 
\text{and} \quad R(\sigma s_i)=R(\sigma)$$  then there exists a unique factorization 
$\sigma s_i=\pi'\rho'$ where either : 
\begin{enumerate}
\item
$\pi=\pi'$ while $\rho^{-1}\rho'\in S$ and $R(\rho)=R(\rho')$ or
\item
$\rho=\rho'$ while $\pi^{-1}\pi'\in S$ and $R(\pi)=R(\pi')$ .
\end{enumerate}
\end{proposition}

\begin{proof}
There exist  two factorizations $\sigma s_i=\pi'\rho'$ such that either $\pi=\pi'$ or $\rho=\rho'$, namely

\begin{itemize}
\item
$\pi'=\pi$, $\rho'=\rho s_i$
\item
 $\pi'=\pi(\rho s_i\rho^{-1})$, $\rho'=\rho$
 \end{itemize}

Suppose that $R(\rho)=R(\rho s_i)$, by Lemma \ref{>2} this happens if and only if $|\rho_i-\rho_{i+1}|\geq 2$, then the factorization
$\sigma s_i=\pi\rho'=\pi(\rho s_i)$ satisfies (1). In the other factorization, $\sigma s_i=\pi'\rho$, one has $\pi'=\pi t$ where $t$ is the transposition $(\rho_i\,\rho_{i+1})$. Since $t$ is not a simple transposition one has $\pi^{-1}\pi'\notin S$ therefore  case $(1)$ of the Proposition holds but not $(2)$.

Suppose now that $R(\rho)\ne R(\rho s_i)$,
this happens if and only if $|\rho_i-\rho_{i+1}|=1$. In this case let
$j=\min (\rho_i,\rho_{i+1})$ and $j+1=\max (\rho_i,\rho_{i+1})$ then
$\rho s_i\rho^{-1}=s_j$  and $|\pi_j-\pi_{j+1}|=|\sigma_i-\sigma_{i+1}|\geq 2$ therefore, by Lemma \ref{>2},
$R(\pi)=R(\pi s_j)$. Now case $(2)$ holds but not $(1)$.

Finally we see that  either $(1)$ or $(2)$ holds and  Proposition \ref{recoil} is proved.
\end{proof}

\end{document}
